\chapter{Related Work}
\label{ch:related-work}

Three lines of research lead to the questions of this thesis: the privacy
harms of health tracking, the privacy communication of health apps, and the
way privacy texts word risk and protection. Each section ends with the
question that its line of research leaves open.

\section{Harms of Health Tracking}
\label{sec:rw-harms}
What others can learn from tracker data is well documented. Salehzadeh Niksirat
et al.\ survey the research on wearable activity trackers and collect the
inferences and threats that their data
allow~\parencite{niksirat2024wearable}. Health apps also pass data on: 33 of 36
depression and smoking-cessation apps sent data to third
parties~\parencite{huckvale2019assessment}.

This evidence comes from measurements and attacks. It has not been collected
into a list of the harms that apply to a health tracker, and it has not been
compared with what providers tell their users. This leads to RQ1.

\section{Privacy Communication of Health Apps}
\label{sec:rw-communication}
Most studies of what health apps communicate look at the privacy policy and
compare it with what the app does. Of 28 apps that sent data to Google, only
12 disclosed this in their policy~\parencite{huckvale2019assessment}. Malki et
al.\ combined the policies and Data Safety sections of 20 female health apps
with an inspection of the apps. They found inconsistencies between policies
and app features, flawed consent and deletion mechanisms and hidden collection
of sensitive data~\parencite{malki2024exploring}. Statements about
consequences appear only in passing: 13 of 19 policies said that law
enforcement or other authorities might access user data, but explained little
about when. In one app, the policy allowed such access while the privacy FAQ
inside the app said that data would never be shared with the government or law
enforcement~\parencite[\S4.2.5]{malki2024exploring}.

In these studies, harms are part of the researchers' threat model, not
something the provider's text is examined for. The contradiction found by
Malki et al.\ also shows that the policy and the interface of the same app can
tell different stories, but the studies reviewed here do not compare the two
surfaces for what they say about harms. This leads to RQ1 and RQ4.

\section{Wording of Risk and Protection}
\label{sec:rw-wording}
A text can disclose a practice and still blur what follows from it. Zou et
al.\ analysed 161 data breach notifications: 112 used hedge words such as
\emph{maybe} or \emph{likely}, and only 9 linked the exposed data to concrete
misuse~\parencite[\S5]{zou2019might}. Protection claims raise a similar
problem, because a PET can work as described and still leave the harms of the
core function possible (Section~\ref{sec:practices-harms}).

For health trackers, the wording of harms has not been studied, and no study
relates a claimed protection to the pathways to harm that it would have to
cut. Hedging as found by Zou et al.\ leads to RQ2, which asks how providers
qualify or downplay the harms they name. Harm trees make it possible to check
which pathways a claimed protection interrupts, which is RQ3.

\section{Research Questions}
\label{sec:rqs}
Together, the three lines leave one question open: \emph{What types of
downstream privacy harms do health tracker providers acknowledge or downplay
when they describe their privacy practices or PETs, and how do they frame the
protection these technologies offer?} Four sub-questions break it down.

\begin{description}[leftmargin=1.5em,itemsep=1pt,topsep=3pt,
                    labelindent=0pt,font=\normalfont\bfseries]
  \item[RQ1] \textit{Coverage.} Which data flows, capabilities and harms do the
    privacy policy and the interface of an app mention? Which of the harms
    identified for the app remain unmentioned?
  \item[RQ2] \textit{Stance.} When the policy or the interface names a harm, how
    do they word it: do they acknowledge, qualify, minimise or deny it, or
    present it as the user's responsibility?
  \item[RQ3] \textit{Protection.} Which protective measures and benefits do the
    policy and the interface claim, and how do the scope and strength of each
    claim relate to the pathways to harm that the measure would have to cut?
  \item[RQ4] \textit{Surfaces.} How do coverage, stance and protection claims
    differ between the privacy policy and the interface of the same app?
\end{description}
